\rtmtight
\begin{tblr}{width=\textwidth,colspec={|Q[c,m,2.1cm]|X[l,m]|},hlines,vlines,hspan=minimal,rowsep=1.5pt,colsep=3pt,row{1}={bg=headgray,font=\bfseries}}
\SetCell{c,m}\hfil\logo\hfil & \textbf{POLITEKNIK NEGERI MALANG}\par\textbf{JURUSAN TEKNOLOGI INFORMASI}\par\textbf{PROGRAM STUDI: D4 TEKNIK INFORMATIKA} \\
\SetCell[c=2]{l} \textbf{RENCANA TUGAS MAHASISWA (RTM) DAN RUBRIK PENILAIAN} & \\
Nama Mata Kuliah & Pemrograman Mobile \\
Kode Mata Kuliah & RTI255003 \quad \textbf{sks / Jam perminggu:} 3 SKS/6 jam \quad \textbf{Semester:} 5 \\
Dosen Pengampu & \begin{enumerate}\item Luqman Affandi, S.Kom., MMSI.\item Dian Hanifudin Subhi, S.Kom., M.Kom.\item Agung Nugroho Pramudhita, S.T., M.T.\item Ade Ismail, S.Kom., M.TI.\end{enumerate} \\
\end{tblr}
\assessmentblock{Tugas Jobsheet/Codelab Flutter}{Hands-on Practice}{SCPMK0209-03901, SCPMK0603-03902}{Mahasiswa menyelesaikan tugas jobsheet/codelab flutter sesuai Sub-CPMK dan konteks mata kuliah.}{Menganalisis kebutuhan, mengerjakan artefak praktik/proyek, menguji hasil, mendokumentasikan evidence, dan mempresentasikan keputusan bila diminta.}{Laporan, artefak digital, repository/prototype, evidence pengujian, dan/atau bahan presentasi.}{Ketepatan penyelesaian codelab Flutter & 12 & Widget, layout, navigation, asset, dan struktur project mengikuti jobsheet dengan hasil berjalan. & Codelab selesai sebagian besar, tetapi ada widget, layout, atau navigation yang belum tepat. & Codelab tidak selesai atau aplikasi tidak dapat dijalankan. \\ \\
Kualitas kode Dart dan komponen UI & 10.5 & Kode Dart rapi, komponen reusable, state sederhana tepat, dan UI responsif pada perangkat target. & Kode berjalan, tetapi komponen masih bercampur atau responsivitas terbatas. & Kode sulit dibaca atau UI tidak sesuai target mobile. \\ \\
Evidence run dan refleksi perbaikan & 7.5 & Screenshot/emulator log, commit, kendala, dan solusi perbaikan terdokumentasi jelas. & Evidence tersedia, tetapi refleksi kendala atau solusi masih umum. & Tidak ada evidence run atau tidak menjelaskan proses penyelesaian. \\}

\assessmentblock{UTS UI/UX Aplikasi Mobile}{UTS praktik}{SCPMK1006-03903}{Mahasiswa menyelesaikan uts ui/ux aplikasi mobile sesuai Sub-CPMK dan konteks mata kuliah.}{Menganalisis kebutuhan, mengerjakan artefak praktik/proyek, menguji hasil, mendokumentasikan evidence, dan mempresentasikan keputusan bila diminta.}{Laporan, artefak digital, repository/prototype, evidence pengujian, dan/atau bahan presentasi.}{Kesesuaian rancangan mobile dengan kebutuhan & 6 & User flow, screen utama, navigation pattern, dan prioritas fitur sesuai konteks aplikasi mobile. & Rancangan cukup sesuai, tetapi flow atau prioritas fitur belum lengkap. & Rancangan tidak sesuai kebutuhan mobile atau flow tidak jelas. \\ \\
Kualitas prototype UI mobile & 5.3 & Layout, hierarchy, touch target, component state, dan konsistensi visual mendukung penggunaan di layar mobile. & Prototype cukup baik, tetapi hierarchy, state, atau konsistensi visual masih kurang. & Prototype sulit digunakan atau tidak memperhatikan karakteristik mobile. \\ \\
Validasi dan presentasi keputusan desain & 3.7 & Rancangan diuji atau direview, temuan ditindaklanjuti, dan keputusan desain dijelaskan dengan evidence. & Review dilakukan, tetapi tindak lanjut atau evidence masih terbatas. & Tidak ada validasi/review atau tidak mampu menjelaskan keputusan desain. \\}

\assessmentblock{PBL Aplikasi Mobile Flutter}{Project Based Learning}{SCPMK0603-03902, SCPMK1006-03903, SCPMK0603-03904}{Mahasiswa menyelesaikan pbl aplikasi mobile flutter sesuai Sub-CPMK dan konteks mata kuliah.}{Menganalisis kebutuhan, mengerjakan artefak praktik/proyek, menguji hasil, mendokumentasikan evidence, dan mempresentasikan keputusan bila diminta.}{Laporan, artefak digital, repository/prototype, evidence pengujian, dan/atau bahan presentasi.}{Kesesuaian fitur aplikasi mobile & 22 & Fitur, user flow, data, dan prioritas aplikasi Flutter sesuai kebutuhan pengguna dan scope proyek. & Fitur utama sesuai, tetapi beberapa flow atau prioritas data belum lengkap. & Fitur tidak menjawab kebutuhan pengguna atau scope proyek tidak jelas. \\ \\
Integrasi Flutter, state, API, dan storage & 19.3 & Widget, navigation, state management, API/local storage, error/loading state, dan native capability terintegrasi stabil. & Integrasi berjalan, tetapi state, error handling, API, atau storage masih terbatas. & Komponen utama tidak terintegrasi atau aplikasi sering gagal. \\ \\
Testing, build, dan demonstrasi aplikasi & 13.7 & Aplikasi diuji di emulator/perangkat, build tersedia, evidence lengkap, dan demo menjelaskan keputusan teknis. & Demo berjalan, tetapi testing/build evidence atau penjelasan teknis belum lengkap. & Aplikasi tidak dapat didemonstrasikan atau tidak ada evidence pengujian. \\}

\begin{tblr}{width=\textwidth,colspec={|Q[l,m,3.2cm]|X[l,m]|},hlines,vlines,hspan=minimal,row{1-3}={bg=headgray},rowsep=1.5pt,colsep=3pt}
Jadwal Pelaksanaan: & Minggu 1--17 sesuai RPS. Setiap evaluasi memiliki RTM dan rubrik multi-kriteria. \\
Lain-Lain Yang Diperlukan: & LMS, repositori/prototype/proyek, perangkat lunak sesuai mata kuliah, dan perangkat presentasi. \\
Pustaka: & Mengacu pustaka utama dan pendukung pada bagian RPS. \\
\end{tblr}
